\section{Introduction and Related Work}

Technical chemistry search rarely respects a single language boundary: a chemist
may query in English, German, French, or Spanish while relevant evidence resides
in a multilingual patent publication, a regional filing, or a document whose
language differs from the query. This is not only a translation problem:
chemistry retrieval must also distinguish a target compound from chemically
related but incorrect alternatives. A retriever for a high-trust workflow such as
patent search or RAG therefore has to do two hard things at once: cross the
language barrier and resist chemical confusability.

Dense retrieval with neural sentence and passage embeddings
\citep{reimers2019sentencebert,karpukhin2020dpr} is attractive in this setting
because multilingual encoders promise a single retrieval index across languages
\citep{labse2022,bgem3_2024}. However, recent work on cross-lingual retrieval
and multilingual RAG shows that retrieval behavior can depend strongly on
language alignment, translation effects, and same-language preference
\citep{whatdrivesclir2025,crosslingualcost2025,bordirlines2024,xrag2025,nepotism2025}.
Aggregate Recall@$k$ can therefore overstate readiness when strong performance
on easier same-language matches masks failures to retrieve cross-language
evidence. It also does not reveal whether chemically related but incorrect
documents outrank the intended evidence.

Existing benchmark lines cover important parts of this space. General embedding
and multilingual retrieval benchmarks such as MTEB, MMTEB, MIRACL, NeuCLIR, and
CLIRMatrix enable systematic comparison of retrieval and embedding models across
languages
\citep{miracl2023,neuclir2023,muennighoff2023mteb,mmteb2025,clirmatrix2020}.
Chemistry-specific resources such as ChemTEB, ChemLit-QA, ChemComp,
ChemKGMultiHopQA, and ChEmbed extend evaluation to chemical embeddings, RAG, and
scientific reasoning
\citep{chemteb2024,chemlitqa2024,khodadad2026chemcomp,astaraki2026iterativerag,chembed2025}.
What remains missing is their combination: multilingual chemistry retrieval data
with related technical evidence across languages, explicit same-language and
cross-language relevance judgments, and chemically plausible hard negatives.
Chemical ontology and alias resources such as ChEBI \citep{chebi2016} provide
the structure needed to define such hard negatives. This use is complementary to
ontology-grounded extraction work such as CEAR \citep{cear2024}, which extracts
chemical entities and roles rather than testing retrieval against chemically
related look-alikes.

Patents are a practical substrate for this evaluation. Patent retrieval has a
long evaluation history, including NTCIR and CLEF-IP
\citep{prime2002,clefip2013}, and recent work has revisited patent
representation learning and embedding evaluation
\citep{paecter2024,patentembeddings2026}. Because patent publications often
contain related technical disclosures across languages, they provide comparable
multilingual technical evidence beyond prior-art search alone.

These gaps motivate X-ChemIR, a patent-grounded benchmark suite for
cross-lingual chemistry retrieval. Our contributions are: (1) two public QAC
retrieval releases derived from Google Patents and EPO data, with
query/document-language metadata, plus an auxiliary ontology-based
chemical-confusability stress test; (2) a reproducible LLM-assisted QAC
generation and human-audit pipeline; (3) language-aware robustness diagnostics
reported next to recall, including same- versus cross-language relevance,
coverage depth, ranking consistency, reading cost, reranker recoverability, and
alignment-only failures; and (4) an evaluation of eight off-the-shelf
multilingual embedding models that quantifies cross-lingual degradation and
chemical-confusion failures.
